(1) 设束缚能 \(E=-\epsilon\)，其中 \(0<\epsilon<V_0\)，阱内动能为 \(V_0-\epsilon\)，故 \(k=\sqrt{2m(V_0-\epsilon)}/\hbar\)；阱外类禁区衰减率 \(\kappa=\sqrt{2m\epsilon}/\hbar\)。两者满足 \(k^2+\kappa^2=2mV_0/\hbar^2\)，是量子态存在的约束。
(2) 偶态内部可取 \(\psi=A\cos(kx)\)，阱外按 \(e^{-\kappa(|x|-a)}\) 衰减。在 \(x=a\) 处连续并匹配一阶导得 \(k\tan(ka)=\kappa\)。奇态内部取 \(\psi=A\sin(kx)\)，同理得 \(-k\cot(ka)=\kappa\)，两套条件必须分别求根，不能把余弦与正弦错置。
(3) 设 \(z=ka\)、\(w=\kappa a=\sqrt{z_0^2-z^2}\)，偶态满足 \(z\tan z=w\)，在首个正切支内总存在交点，因此任意 \(z_0>0\) 都至少有一个偶束缚态。首个奇态在束缚能从零开始出现时 \(w\to0\)，对应 \(z=\pi/2\)，故阈值为 \(z_0>\pi/2\) 才有严格负能奇态，等号是非平方可积的阈值态。
(4) \(V_0\to\infty\) 且井宽 \(2a\) 固定时，阱外衰减加强，能级相对井底趋近无限深阱的 \((n\pi\hbar)^2/(8ma^2)\)；但对有限 \(V_0\)，阱外概率密度非零且各能级是连续匹配方程的根，因此不能直接套用无限深势阱的精确能谱。
